%% Unnumbered "Ethical Considerations" section. *ACL exempts it from the page
%% limit when placed after the conclusion, so it is inputted only by
%% main_acl.tex; the AAAI build has no such exemption.
\textbf{Deployment risk is the motivation and also the hazard.} The agents we study feed
financial analytics and investment decisions, and our central result is that they commit
confidently to an unintended reading of an ambiguous question far more often than they ask.
A wrong number produced this way is well grounded in a real filing, so it is unusually hard
for a downstream user to catch. We therefore frame the benchmark as a reason to prefer
systems that ask and abstain over systems that answer, and we caution that a high score on
\fininteract measures a narrow clarification ability, not readiness for deployment in a
setting where a confident wrong figure carries financial consequences.

\textbf{Intended use and misuse.} \fininteract is an evaluation resource, not a training
corpus and not financial advice, as stated in its datasheet (Appendix~\ref{app:datasheet}).
Because every instance is ambiguous by construction, the benchmark does not measure a
calibrated decision of \emph{when} to ask, so it should not be used to certify that a system
asks appropriately in general deployment; we make this limitation explicit in
Appendix~\ref{app:limitations}.

\textbf{Data provenance and privacy.} All source documents are public regulatory filings:
SEC EDGAR filings for the English split and CSRC annual reports retrieved via CNINFO and
akshare for the Chinese split. No proprietary, paywalled, or personal data is used. Instances
consist of questions about corporate financial metrics and carry company identifiers rather
than individuals, so the released artifact contains no personally identifying information and
no objectionable content. Third-party corpora used only during construction are not
redistributed.

\textbf{Compute and environmental cost.} The evaluation is API-based rather than
training-heavy. We instrument and report the full construction cost in
Appendix~\ref{app:cost}, at \$0.31 per accepted instance and about \$54 for the released
corpus. The only model training in this work is 4-bit QLoRA on a 4B-parameter policy, run for
tens of optimization steps rather than to convergence, so the associated energy cost is small
relative to a typical pretraining or full fine-tuning study.
